\section{Tangent transfer and matching normal drops}

We record a concrete straightedge-and-compass configuration that makes the affine ``transport'' of the short-time normal departure from the tangent completely explicit.
The purpose of this section is to obtain the geometric local-drop estimate; the kinematic conversion to acceleration then uses \Cref{lem:deflection-identity,lem:bd-over-br2} from Section~3.

\subsection{Circle companions and tangent transfer}

Let $\ell$ be the major axis, and let $\Phi$ map the auxiliary circle $\mathcal{C}$ to the ellipse. From $A$ and $B$ drop perpendiculars to $\ell$ and extend these same vertical lines to meet the auxiliary circle $\mathcal C$ at $A_0$ and $B_0$ (choose the intersections on the same side of $\ell$ as $A$ and $B$). By construction, $\Phi(A_0)=A$ and $\Phi(B_0)=B$.
Draw the tangent to $\mathcal C$ at $A_0$ and let it meet $\ell$ at $T$.

\begin{lemma}[Tangent transfer with fixed $T$]
  \label{lem:tangent-transfer}
  Under $\Phi$, the circle tangent at $A_0$ maps to the ellipse tangent at $A$. Moreover, since $\Phi$ fixes $\ell$ pointwise, the intersection point $T\in\ell$ is the same for both tangents.
\end{lemma}
\begin{proof}
  $TA_0$ is the tangent to the circle at $A_0$. For points $B_0$ on the circle taken sufficiently near $A_0$, all such $B_0$ lie on the same side of the line $TA_0$.

  The affine bijection $\Phi$ sends lines to lines and half-planes to half-planes, so it preserves the relation ``lying on the same side of a line.'' Hence the line $\Phi(T)\Phi(A_0)$ meets the ellipse at $A=\Phi(A_0)$. Since $\Phi(T)=T$ and all nearby image points $B=\Phi(B_0)$ lie on one side of $TA$, the line $TA$ is tangent to the ellipse at $A$.

  Here we implicitly use the fact that a conic is smooth and therefore has a unique and well-defined tangent line at each point.\qedhere
\end{proof}

\subsection{The \texorpdfstring{$B,C,D$}{B,C,D} construction (ellipse and circle)}

Define $C$ as the intersection of the ellipse tangent $TA$ with the line through $B$ parallel to $\ell$, and define $D:=TA\cap FB$. Likewise, on the circle tangent $TA_0$ define $C_0$ as the intersection with the line through $B_0$ parallel to $\ell$, and define $D_0:=TA_0\cap OB_0$. By construction, $BC\parallel\ell$ and $B_0C_0\parallel\ell$.

\subsection{The point \texorpdfstring{$J$}{J} is on \texorpdfstring{$\mathcal{C}$}{C}}

Extend the ray $AF$ beyond $A$ and mark a point $F''$ on this ray such that $AF''=AF'$ (so $\triangle AF'F''$ is isosceles). Let $J:=AT\cap F'F''$.

\begin{lemma}[Ellipse tangent bisects the focal angle]
  The tangent line $AT$ bisects $\angle F''AF'$.
\end{lemma}

In particular (by the isosceles geometry), $AT\perp F'F''$ and $J$ is the midpoint of $F'F''$. But we also know $O$ bisects $FF'$, so $OJ\parallel FA$ and $OJ=FF''/2=a$, which proves that $J$ lies on the circle $\mathcal{C}$.

\subsection{Matching the normal drops}

The triangle similarity relations in the tangent line through $T$ give
\begin{equation}
  B_0D_0 = B_0C_0\,\frac{OD_0}{OT},
  \qquad
  BD = BC\,\frac{FD}{FT}.
\end{equation}
Dividing and using that $\Phi(B_0)=B$ and $\Phi$ maps the line through $B_0$ parallel to $\ell$ to the line through $B$ parallel to $\ell$ (since $\Phi$ fixes $\ell$ and preserves parallelism), while also mapping the tangent line $TA_0$ to $TA$, we have $\Phi(C_0)=C$. Hence the horizontal segment $B_0C_0$ is carried to $BC$, and since $\Phi$ preserves horizontal lengths, $B_0C_0=BC$.
\begin{equation}
  \frac{B_0D_0}{BD}=\frac{OD_0}{OT}\cdot\frac{FT}{FD}.
\end{equation}
Letting $B\to A$ (so $D_0\to A_0$ and $D\to A$), we obtain
\begin{equation}
  \lim_{B\to A}\frac{B_0D_0}{BD}=\frac{OA_0}{OT}\cdot\frac{FT}{FA}.
\end{equation}

\begin{lemma}[Matching drops]
  \label{lem:matching-drops}
  In the infinitesimal-step limit, the circle and ellipse normal drops from the tangent agree:
  \begin{equation}
    \boxed{\;\lim_{B\to A}\frac{B_0D_0}{BD}=1.\;}
  \end{equation}
\end{lemma}
\begin{proof}
  Now $F,O,T\in\ell$, and by construction $J\in AT$ and $OJ\parallel FA$. Hence triangles $\triangle TFA$ and $\triangle TOJ$ are similar. Therefore
  \[
    \frac{FT}{FA}=\frac{OT}{OJ}.
  \]
  Since $A_0,J\in\mathcal C$ we have $OA_0=OJ=a$, and consequently the right side of (5)
  \[
    \frac{OA_0}{OT}\cdot\frac{FT}{FA}=\frac{a}{OT}\cdot\frac{OT}{a}=1.
  \]
  which proves (6).
\end{proof}

\subsection{Transporting swept area from the ellipse to the circle}
\label{sec:transport-and-sagitta}

Let $A$ be the current point on the ellipse and $B$ the position after a small time increment $t$. Let the swept area about the focus be
\begin{equation}
  S = \operatorname{area}(\triangle F A B)=\frac{L}{2}\,t.
\end{equation}

We compare this to a corresponding small triangle on the circle. Let
\begin{equation}
  S_0 := \operatorname{area}(\triangle O A_0 B_0),
\end{equation}
Two geometric scalings relate $S$ and $S_0$:
\begin{enumerate}
  \item Height scaling from $O$ to $F$. In the construction used in our discussion, the line through $F$ parallel to a fixed circle radius implies that the perpendicular height from $F$ to the chord $AB$ is $(r/a)$ times the perpendicular height from $O$ to the same chord direction. Hence
    \begin{equation}
      \operatorname{area}(\triangle F A B)=\frac{r}{a}\,\operatorname{area}(\triangle O A B).
    \end{equation}

  \item Affine area scaling. Since $\Phi$ scales area by $b/a$,
    \begin{equation}
      \operatorname{area}(\triangle O A B)=\frac{b}{a}\,\operatorname{area}(\triangle O A_0 B_0)=\frac{b}{a}\,S_0.
    \end{equation}
\end{enumerate}

Combining,
\begin{equation}
  S=\operatorname{area}(\triangle F A B)=\frac{r}{a}\cdot\frac{b}{a}\,S_0=\frac{rb}{a^2}\,S_0,
\end{equation}
so
\begin{align}
  S_0 &= \frac{a^2}{rb}\,S \label{eq:s0-from-area}
  = \frac{a^2}{rb}\cdot \frac{L}{2}\,t
  = \frac{L a^2}{2 b r}\,t.
\end{align}

\subsection{Circle tangent and sagitta}
\label{subsec:sagitta}

As $B_0\to A_0$, \eqref{eq:s0-from-area} gives
\begin{equation}\label{eq:a0d0}
  A_0D_0 = \frac{2S_0}{a}=
  \frac{2}{a}\cdot\frac{L a^2}{2 b r}\,t
  =\frac{L a}{b r}\,t.
\end{equation}

Now use elementary circle geometry relating the \emph{normal drop} from the circle to its tangent (the sagitta). For a circle of radius $a$, the perpendicular distance from the circle point reached to the tangent line satisfies (to leading order)
\begin{equation}
  B_0D_0 \approx \frac{(A_0D_0)^2}{2a}.
\end{equation}
Substituting \eqref{eq:a0d0} gives
\begin{equation}
  B_0D_0 \approx \frac{1}{2a}\left(\frac{L a}{b r}t\right)^2
  = \frac{L^2 a}{2 b^2 r^2}\,t^2.
\end{equation}

\subsection{Proof 1 (ellipse via affine transported drops)}
\label{subsec:proof1-ellipse-affine}

\begin{proposition}[Proof 1: direct problem for the ellipse]
\label{prop:proof1-ellipse-inverse-square}
Under Hypotheses H1--H2, the acceleration is focus-directed and has inverse-square magnitude:
\begin{equation}
  \mathbf a_F(r) = -\mu\,\frac{\mathbf r}{r^3},
  \qquad
  \mu=\frac{L^2 a}{b^2}=\frac{L^2}{p},
\end{equation}
where $p=b^2/a$ is the semi-latus rectum and $\mathbf r$ is the position vector from $F$.
\end{proposition}
\begin{proof}
From \Cref{lem:matching-drops} and the sagitta estimate above,
\[
  BD\sim B_0D_0\sim \frac{L^2 a}{2 b^2 r^2}\,t^2.
\]
Also, by the area-law relation (equivalently, the first asymptotic in \Cref{lem:bd-over-br2}),
\[
  BR\sim \frac{L}{r}\,t.
\]
Therefore
\[
  \frac{BD}{BR^2}\to \frac{a}{2b^2}=\frac{1}{2p},
  \qquad \left(p=\frac{b^2}{a}\right).
\]
Applying \Cref{lem:bd-over-br2} with $c=\frac{1}{2p}$ gives
\[
  a_F=\frac{2cL^2}{r^2}
  =\frac{L^2}{p}\cdot\frac{1}{r^2}
  =\frac{L^2 a}{b^2}\cdot\frac{1}{r^2}.
\]
Combining this magnitude with Section~3's central-direction conclusion yields
\[
  \mathbf a_F(r) = -\mu\,\frac{\mathbf r}{r^3},
  \qquad
  \mu=\frac{L^2a}{b^2}=\frac{L^2}{p}.
\]
For the areal constant $k$ with $L=2k$, this is equivalently $\mu=4k^2/p$.
\end{proof}
